\documentclass[11pt]{article}
\usepackage[margin=1in]{geometry}
\usepackage[T1]{fontenc}
\usepackage{enumitem}
\setlength{\parindent}{0pt}
\setlength{\parskip}{6pt}

\title{LibraryCard --- Written Questions}
\author{}
\date{}

\begin{document}
\maketitle

\begin{enumerate}[leftmargin=2em, itemsep=1em]

\item \textbf{Which C++ feature did you use to meet B4 (the number never changes)?
Why does that choice force you to set the number in a particular place in the constructor?}

I declared the member as \texttt{const int number} in \texttt{LibraryCard.h}.
A \texttt{const} data member cannot be assigned after the object is created, so
the only chance to give it a value is in the constructor's member-initializer
list (I used \texttt{number(nextNumber++)} there); assigning it in the
constructor body would not compile. This makes the compiler, not my
discipline, enforce that nothing ever changes the number.

\item \textbf{How did you meet B5 (each rule written once)? Name the feature or
features you used.}

Each rule lives in its own \emph{private static member function}:
\texttt{isValidName()} holds the ``at least 2 characters'' rule and
\texttt{cleanFines()} holds the ``negative becomes 0'' rule. The constructor
calls both in its member-initializer list, and \texttt{setHolder()} calls
\texttt{isValidName()}. In addition, \emph{delegating constructors} route the
one-argument and default constructors through the two-argument constructor,
so all three creation paths share a single body. No check is copy-pasted.

\item \textbf{The test program's output ends with five closed lines. Explain why
they come out in that order, and why Card 5004 closed appears much earlier.}

Automatic (local) objects are destroyed in \emph{reverse order of construction}
when the block they live in ends. \texttt{temp} (card 5004) is declared inside
its own inner block, so its destructor runs at that block's closing brace,
which is why ``Card 5004 closed'' appears right after the second
\texttt{Active: 4} line. The five locals in \texttt{main} were constructed in
the order ana (5001), ben (5002), x (5003), \texttt{lobby[0]} (5005),
\texttt{lobby[1]} (5006), so they are destroyed in exactly the reverse order:
\texttt{lobby} first, array elements in reverse index order (5006, 5005),
then x (5003), ben (5002), and ana (5001).

\item \textbf{If printReport failed to compile at some point, what was the error
and what did you change? If it never failed, what did you do in advance so it
wouldn't?}

It never failed. In advance, I marked every member function that only reads the
card --- the four getters and \texttt{print()} --- as \texttt{const} in the
class declaration. The parameter of \texttt{printReport} is
\texttt{const LibraryCard\&}, and a const object may only call const member
functions, so without those \texttt{const} qualifiers the call
\texttt{card.print()} would not have compiled.

\item \textbf{If copying were allowed, what would go wrong with the card numbers?
Use the words copy constructor in your answer.}

The compiler-generated copy constructor would copy the \texttt{number} field
verbatim into the new object, so two live card objects would report the same
card number. That breaks the uniqueness rule (B3): after
\texttt{LibraryCard copy = ana;} both cards would claim to be card 5001, and
there would be no way to tell which physical card a record belongs to.

\item \textbf{Your class has both constructors and a setter for the holder's name.
Why isn't setHolder enough on its own?}

A member function can only run on an object that already exists, so the object
must be created in a valid state before \texttt{setHolder} could ever be called:
creation must sanitize the name, assign the card number, and zero the book
count without any chance of failure (B2), whereas \texttt{setHolder} is allowed
to reject a name and return \texttt{false}. The card number makes this
impossible anyway, because it is a \texttt{const} member that no setter can
ever assign --- only a constructor's initializer list can give it its value.

\item \textbf{Suppose you forgot to add LibraryCard.cpp to add\_executable. Which
tool reports the problem (compiler or linker), and what does the error say?}

The \emph{linker} reports it. \texttt{main.cpp} still compiles, because the
declarations in \texttt{LibraryCard.h} are enough for the compiler; but with no
definitions anywhere, linking fails with errors such as:
\begin{verbatim}
undefined reference to `LibraryCard::LibraryCard(std::__cxx11::basic_string<char, std::char_traits<char>, std::allocator<char> >)'
undefined reference to `LibraryCard::activeCards()'
undefined reference to `LibraryCard::count'
collect2: error: ld returned 1 exit status
\end{verbatim}

\end{enumerate}

\section*{E1 --- First error line for each must-not-compile check}

Each line was uncommented one at a time and built; the first error line from
each build is copied below (line numbers refer to my \texttt{main.cpp}).

\begin{enumerate}[leftmargin=2em]
\item \texttt{LibraryCard copy = ana;}
\begin{verbatim}
main.cpp:58:24: error: use of deleted function 'LibraryCard::LibraryCard(const LibraryCard&)'
\end{verbatim}

\item \texttt{lobby[0] = ana;}
\begin{verbatim}
main.cpp:59:16: error: use of deleted function 'LibraryCard& LibraryCard::operator=(const LibraryCard&)'
\end{verbatim}

\item \texttt{std::string name = "Zoe";  printReport(name);}
\begin{verbatim}
main.cpp:60:44: error: invalid initialization of reference of type 'const LibraryCard&'
from expression of type 'std::string' {aka 'std::__cxx11::basic_string<char>'}
\end{verbatim}
\end{enumerate}

The third error occurs because the \texttt{LibraryCard(std::string)}
constructor is \texttt{explicit}, so a \texttt{std::string} cannot implicitly
convert to a \texttt{LibraryCard} when initializing the
\texttt{const LibraryCard\&} parameter (B6).

\end{document}
